% ---------------------------------------------------------------------------
%  QML Lab research project 06: Quantum Autoencoder
%  Build:  latexmk -pdf main.tex   (needs qmlproject.sty, harvardqml.bst, references.bib)
% ---------------------------------------------------------------------------
\documentclass[11pt,a4paper]{article}
\usepackage{qmlproject}
\qmlheader{Project 06: Quantum Autoencoder}

\hypersetup{pdftitle={QML Lab project 06: Quantum Autoencoder}}

\begin{document}

\begin{center}
  {\small\textcolor{black!55}{QML Lab \textbullet\ research project 06}}\\[0.5em]
  {\LARGE\bfseries Quantum Autoencoder}\\[0.6em]
  {\large Compressing quantum states into fewer qubits and recovering them.}\\[0.6em]
  {\small Living document \textbullet\ last updated \today}
\end{center}

\begin{projectcard}
\small
\begin{tabularx}{\linewidth}{@{}>{\bfseries}l X >{\bfseries}l X@{}}
Quadrant  & \QQ\ (quantum data, quantum processor) & Track & Variational \\
Difficulty & Medium & Tools & PennyLane \\
Builds on & Modules 2 and 5 of the QML Lab & Workspace & \path{qml_projects/quantum-autoencoder/} \\
\end{tabularx}
\end{projectcard}

\paragraph{Summary.}
A quantum autoencoder learns a unitary that squeezes a family of quantum states into fewer qubits,
so that the discarded qubits always end up in a fixed reference state. Because the inputs and the
processing are both quantum, this is a \QQ\ project. It studies which state families compress, how
far, and how robust the trained encoder is to noise.

% ===========================================================================
\section{Background and key concepts}
% ===========================================================================
\subsection{Compression of quantum data}
Let $\{\ket{\psi_i}\}$ be a family of $n$-qubit states, for example ground states of a Hamiltonian at
different couplings. If all of them lie in a subspace of dimension at most $2^{n_A}$, a unitary $U$ can
map them onto $n_A$ ``latent'' qubits, leaving $n_B=n-n_A$ ``trash'' qubits in $\ket{0}$. Decoding
applies $U^\dagger$ after resetting the trash qubits \citep{romero2017}.

\subsection{Training objective}
For a parameterised encoder $U(\theta)$, the trash fidelity of state $\psi_i$ is
\begin{equation}
  F_i(\theta)=\Tr\!\Big[\big(I_A\otimes\ket{0}\!\bra{0}_B\big)\,U(\theta)\ket{\psi_i}\!\bra{\psi_i}U(\theta)^\dagger\Big],
\end{equation}
and training minimises $C(\theta)=1-\frac1N\sum_iF_i(\theta)$. If $C=0$, decoding reproduces every state
exactly. On hardware, overlaps such as this are often measured with the swap test
\citep{buhrman2001}, which returns 0 with probability $\tfrac12\big(1+|\braket{a|b}|^2\big)$.
\citet{pepper2019} realised a quantum autoencoder experimentally with photonic qutrits.

\subsection{Why this is QQ}
The encoder acts on quantum states directly, without measuring them. In simulation the input states
are prepared classically, which hides the real use case: compressing the output of a quantum
simulation or a quantum sensor before storing or transmitting it.

% ===========================================================================
\section{Research questions}
% ===========================================================================
The project is open-ended: these questions are the starting point, not the finish line. Add new
ones to the README and the journal as they come up.

\begin{rqlist}
  \item Which state families compress well (Ising ground states, GHZ-like, W-like, Haar-random)?
  \item How does reconstruction fidelity fall when compressing further (3 to 2 versus 3 to 1 qubits)?
  \item How sensitive is a trained encoder to depolarizing noise?
\end{rqlist}

% ===========================================================================
\section{Method}
% ===========================================================================
\subsection{Data}
State families on 3--4 qubits: ground states of the transverse-field Ising model for a range of field
strengths (shared with Project 12), GHZ states with varying relative phase, W-like states, and
Haar-random states as a negative control. Separate training and test states from each family.

\subsection{Models}
Hardware-efficient encoder with $L$ layers of $R_Y$/$R_Z$ rotations and CNOT rings on all qubits;
cost from exact trash-qubit probabilities, and later from finite-shot swap tests; Adam optimiser.

\subsection{Baselines}
Classical reference: the rank of the state family's density matrix, which bounds the best possible
compression; Haar-random states, which should not compress.

\subsection{Metrics}
\begin{itemize}
  \item Trash fidelity and reconstruction fidelity on test states.
  \item Compression ratio $n_A/n$.
  \item Fidelity under depolarizing noise $p$.
\end{itemize}

% ===========================================================================
\section{Planned experiments}
% ===========================================================================
Each experiment becomes one \texttt{[[experiments]]} entry in the workspace, so its runs, parameters
and metrics are stored and compared in the portal.

\begin{table}[h]
\centering\small
\begin{tabularx}{\linewidth}{@{}l>{\raggedright}X>{\raggedright}X>{\raggedright\arraybackslash}X@{}}
\toprule
ID & Experiment & What varies & Main metrics \\
\midrule
  E1 & 3 to 2 qubits & Ising family & test reconstruction fidelity \\
  E2 & 3 to 1 qubit & Ising and GHZ families & fidelity drop \\
  E3 & Negative control & Haar-random states & fidelity vs rank bound \\
  E4 & Noise & depolarizing $p\in[0,0.05]$ & fidelity \\
  E5 & Swap-test cost & shots $M$ & training stability \\
\bottomrule
\end{tabularx}
\caption{Experiment plan. Results go into the journal and the results log (\cref{sec:log}).}
\end{table}

% ===========================================================================
\section{Running it in the QML Lab}
% ===========================================================================
The workspace lives in \path{qml_projects/quantum-autoencoder/}. Scripts import three helpers from
\path{qml_lab/plots.py}:
\begin{itemize}
  \item \texttt{params()} reads the parameters from the portal's run form (defaults when run by hand);
  \item \texttt{save()} stores figures with the run;
  \item \texttt{log\_metrics()} records numbers that appear in the run table.
\end{itemize}

\begin{lstlisting}
qml_projects/quantum-autoencoder/
  project.toml          metadata + experiments
  README.md             goal, open questions, ideas
  experiments/          runnable scripts
  docs/main.tex         this document
\end{lstlisting}

First experiment: This experiment is \emph{planned}: write the script, add the block below to \texttt{project.toml}, then run \texttt{python manage.py sync\_projects}.

\begin{lstlisting}[language={}]
[[experiments]]
slug = "compress-ising"
title = "Compress Ising ground states from 3 to 2 qubits"
script = "experiments/e01_compress_ising.py"
figure = "fidelity.png"

[experiments.params]
n_qubits = 3
n_trash = 1
n_layers = 3
steps = 200
family = "ising"
noise = 0.0
seed = 0
\end{lstlisting}

\begin{lstlisting}[language=bash]
# from the portal folder: pick up changes to project.toml
python manage.py sync_projects
cd qml_projects/quantum-autoencoder
# run by hand with the default parameters
PYTHONPATH=../../qml_lab python experiments/e01_compress_ising.py
\end{lstlisting}

% ===========================================================================
\section{Quality bar}
% ===========================================================================
\begin{itemize}
  \item A classical baseline with a similar number of trainable parameters, tuned with the same care \citep{bowles2024}.
  \item At least five random seeds per configuration, reported as mean $\pm$ standard deviation.
  \item How results change with the number of qubits.
  \item At least one run under a realistic noise model.
  \item Always compare achieved fidelity with the rank bound of the state family.
\end{itemize}

% ===========================================================================
\section{Open directions}
% ===========================================================================
\begin{itemize}
  \item Denoising autoencoders that remove noise from quantum states.
  \item Compress Ising states before feeding them to a QCNN (Project 12).
  \item Run a small encoder on hardware with a swap-test cost.
\end{itemize}

% ===========================================================================
\section{Risks and limitations}
% ===========================================================================
\begin{itemize}
  \item Simulation prepares inputs classically, so the practical advantage of \QQ\ processing is not visible.
  \item Deeper encoders on more qubits may hit barren plateaus.
\end{itemize}

% ===========================================================================
\section{Results log}
\label{sec:log}
% ===========================================================================
Summarise findings here as they become stable; day-to-day notes belong in the portal journal.

\begin{table}[h]
\centering\small
\begin{tabularx}{\linewidth}{@{}l l >{\raggedright}X >{\raggedright\arraybackslash}X@{}}
\toprule
Date & Run & Finding & Next step \\
\midrule
 & & & \\
 & & & \\
 & & & \\
\bottomrule
\end{tabularx}
\end{table}

\bibliographystyle{harvardqml}
\bibliography{references}

\end{document}
